\section{Related Work}
\label{sec:relatedwork}


\subsection{Audits and Accountability}
\label{sec:relatedwork-audits}
Researchers, law and policymakers, and civil society groups increasingly call for audits as mechanisms for identifying and mitigating risks of algorithmic systems and holding developers and users accountable for the algorithms they develop and use~\cite{raji2022outsider, wyden2023algorithmic}. \citet{birhane2024ai} define an audit as \textit{``any independent assessment of an identified audit target via an evaluation of articulated expectations with the implicit or explicit objective of accountability.''} Audits are not restricted to a specific method or type of audit target~\cite{ojewale2025towards, birhane2024ai, goodman2022algorithmic}. An audit may focus narrowly on a specific model or technical system~\cite{buolamwini2018gender}, or more broadly on the sociotechnical system encompassing technical components and human users~\cite{lam2023sociotechnical}. Audits may assess computational notions of fairness, but may also focus on other criteria such as functionality and legality~\cite{radiya2023sociotechnical, ojewale2025towards, abebe2022adversarial}. Additionally, audits may be conducted by a variety of actors, including an entity within the organization that built the examined algorithmic system, or those outside the organization. Most importantly, an audit does not end with a thorough performance analysis. Instead, evaluation of the audit target is situated within a broader accountability process in which the audited entity or organization faces consequences or is otherwise responsible for acting on the outcomes of the evaluation~\cite{ojewale2025towards}.


However, audits do not always successfully create accountability~\cite{latonero2021human, selbst2021institutional, goodman2022algorithmic, raji2022outsider}. When an audit is poorly designed or executed, a claim that a system has been audited can, at best, be meaningless, and at worst, lead to ``audit washing'' where the audit provides false assurance by lending credibility to a dysfunctional or otherwise inadequate system~\cite{goodman2022algorithmic}. For example, recent studies investigating audits of AI hiring systems performed in accordance with New York City's Local Law 144 (LL144) discuss how audits might be able to comply with the law's audit requirements despite being incomplete assessments of algorithmic bias~\cite{wright2024null, gerchick2025auditing}. Collectively, this body of work calls for careful attention to how audits are designed and conducted, and highlights how inadequate audits can hide and lend credibility to dysfunctional or inadequate AI systems. We build on this work by exploring how the design of an auditing standard can contribute to audit washing by lending credibility to audits that can be interpreted to comply with the standard despite failing to produce the outcomes envisioned by the standard's creators. 


\subsection{Designing Effective Tools for Responsible AI}
A growing body of work in human-computer interaction and computer-supported cooperative work seeks to design tools such as software, frameworks, and checklists to help various stakeholders promote fairness, transparency, and accountability throughout the development, deployment, and use of AI systems~\cite{deng2022exploring, lee2021landscape,madaio2020co,raji2020closing, gebru2021datasheets, mitchell2019model}. Recent work examines these tools in the context of AI auditing, highlighting tools that support various stages of the auditing process such as audit standards that outline best practices, and repositories that facilitate broader communication of audit results~\cite{ojewale2025towards}. Several groups have explicitly highlighted the importance of audit standards as tools for improving audit quality and consistency, thereby minimizing the risk of audit washing~\cite{costanza2022audits, raji2022outsider, goodman2022algorithmic, ojewale2025towards, lam2024framework}.

However, numerous studies have shown that the existence of an RAI tool, alone, does not guarantee its effectiveness. A tool must be adopted and used, often within organizational structures and cultures that can hinder a tool's effectiveness in practice~\cite{wong2023seeing, madaio2020co, rakova2021responsible, metcalf2021algorithmic}. Collectively, this body of work emphasizes the importance of examining how a tool is (or might be) used in practice and designing tools with real-world organizational cultures, settings, and incentives in mind. In parallel and in response to these studies, recent work has called for the human-computer interaction community to expand its work to evaluating the effectiveness, and not just the usability, of these tools~\cite{berman2024scoping}.

\subsubsection{Designing Effective Audit Standards}
 Similarly, the existence of an audit standard, or an audit's compliance with a standard, does not necessarily ensure accountability~\cite{ojewale2025towards}. In fact, poorly designed standards can lead to ineffective audits and can lend further credibility to inadequate audits. Standards defined by the audit target~\cite{goodman2022algorithmic, hoofnagle2016assessing}, and standards with overly broad or vague language~\cite{goodman2022algorithmic, costanza2022audits, raji2022outsider, wright2024null, morrison2020vacuous}, can fail to ensure adequate audits and can add further credibility to an inadequate system. For example, \citet{wright2024null} examine audit requirements defined in New York City's Local Law 144 and discuss how the language grants companies substantial discretion over whether their technology falls within the scope of the law's audit requirements. This discretion can undermine accountability by allowing companies to decide that a technology is out of scope, even if the system is potentially biased and meaningfully impacts employment outcomes.
 

To ensure that auditing standards lead to high quality audits and provide appropriate assurance about compliant audits, groups have studied the outcomes of requirements for audits in existing laws and synthesized lessons from other auditing domains such as finance and medicine~\cite{raji2022outsider, wright2024null, groves2024auditing, lam2024framework, gerchick2025auditing}. This body of work has collectively identified several key considerations for the design of effective standards. Some groups have called for standards that better account for their users' needs and work contexts~\cite{ojewale2025towards, schor2024mind}. Others have discussed how the design of audit requirements must enable auditors to draw unambiguous conclusions about the audit target and enable audit stakeholders to draw unambiguous conclusions about a completed audit's compliance with the standard. Drawing on their experiences conducting audits under LL144, \citet{lam2024framework} emphasize that requirements must be ``verifiable or observable conditions'' that ``must enable auditors to form an unambiguous opinion about whether a given criterion is satisfied.'' \citet{wright2024null} further emphasizes the importance of unambiguously verifiable criteria through their discussion of how the discretion that the language in LL144's requirements grants to companies makes it difficult for outsiders to evaluate whether the company complies with the law, a determination that drives critical actions such as regulator demands that a company stop using a given technology. Taken together, this line of work highlights the importance of verifiable and observable requirements that support unambiguous conclusions \cite{lam2024framework, wright2024null, raji2022outsider}, requirements that collectively support holistic evaluations \cite{lam2024framework}, and context-specific requirements that are compatible with the practices and contexts of those using the standard \cite{schor2024mind, ojewale2025towards}. We seek to expand these efforts to build an evidence base for the design of auditing standards by studying aspects of the design of auditing standards that undermine its effectiveness. 



